\section{One exact search over overlapping covers}\label{sec:implementation}

The localization theorem makes the old spatial search complete for descent,
but independent restarts still repeat substantial work. A source with many
small overlapping covers can revisit the same marked trace prefix once for
each cover containing its footprint. The new implementation represents the
union of those finite search languages directly.

\subsection{The cover intersection invariant}

Let $\mathcal R$ be the complete family of permitted original regions,
including the empty region. For each original tetrahedron $v$, build the
inverted incidence list
\[
 I(v)=\{R\in\mathcal R:v\in R\}.
\]
A search frame with actual footprint $F$ retains precisely
\begin{equation}\label{eq:cover-intersection}
 C(F)=\{R\in\mathcal R:F\subseteq R\}
     =\bigcap_{v\in F} I(v).
\end{equation}
The empty intersection means all covers. The implementation represents that
root case by a sentinel instead of copying the complete index. If an event
first consumes the original cells in $S$, its compatible covers become
\[
 C(F\cup S)=C(F)\cap\bigcap_{v\in S}I(v).
\]
An empty result rejects that branch. Consuming a born cell adds no new
original identity: its ancestors already entered the footprint at their
own first consumption. Intersections begin with the smallest available
incidence set, reducing intermediate allocations without changing the set.

\begin{proposition}[Exact union language]\label{prop:union}
The indexed cover search permits exactly the traces admitted by at least
one independently restarted search over $\mathcal R$. It preserves the
relative commutation equivalence used by the sleep-set search.
\end{proposition}
\begin{proof}
Induction on the trace length proves \eqref{eq:cover-intersection}. A trace
is admitted exactly when some region contains every consumed original cell,
which is the definition of membership in the union of the restarted
languages. Every prefix consumes a subset of that final footprint, so it
is admitted under the same covering region.

For two adjacent independent events in an admitted trace, either order has
the same final footprint, while the two intermediate footprints are subsets
of it. The same cover therefore admits both orders. Their geometry and
markings agree by \cref{lem:forward}. Thus all swaps required by the
sleep-set argument remain within the union language.
\end{proof}

Two events may be individually allowed at a prefix although their combined
footprint fits no cover. The proposition does not assert that every such
pair can be combined. It asserts commutation when the whole two-event trace
is admitted, which is exactly the property the reduction needs.

An explicit index has at most $\sum_{R\in\mathcal R}|R|$ incidences. For
connected radius $R_0=O(U)$ this is $t2^{O(U)}$ times a polynomial parameter
factor. Intersecting finite sets and copying compatible-cover states keeps
the same fixed-parameter time form, although it can consume more memory
than streaming one region at a time. The theorem does not promise that a
fully materialized index is the best implementation for every input.

There is also a sharper accounting for cover intersections. Let $C_f$ be
the compatible-cover set at visited frame $f$, and let $K_R$ bound the
commuting trace classes admitted by a fixed cover $R$. Then
\[
 \sum_f |C_f|
 =\sum_{R\in\mathcal R}\#\{f:R\in C_f\}
 \le\sum_{R\in\mathcal R}K_R.
\]
For a fixed cover, the global sleep search still visits at most one
representative of each admitted class. An intersection can scan the smaller
set; ordered incidence tables with binary search give a deterministic
logarithmic overhead per tested membership. Thus incidence work can be
charged to compatible cover/frame pairs instead of multiplying every
frame by the size of the entire region family. The prototype uses exact
Python sets, with equality resolving hash collisions.

\subsection{Lazy transitions and endpoint observations}

Each queued branch stores its parent state and the proposed geometric event.
The native replacement and integral cochain transport are performed only
when that branch is visited. A positive early answer can therefore avoid
constructing all unused siblings. The node limit is checked before producing
the next unvisited child. This matters on the upward-required fixtures:
the source has many upward choices, while a successful trace uses only
three events.

Exact marked endpoint keys suppress repeated endpoint queries. They retain
the source incarnation identities and transported cochain, using exact
interned structures rather than probabilistic hash equality. The search
does \emph{not} delete continuation frames merely because an endpoint query
was already made. Remaining upward budget, sleep obligations, and compatible
covers can differ between histories reaching similar geometry. Proving a
safe stronger quotient is a separate problem.

A complete conclusion for a caller's predicate requires that predicate to
be invariant under the represented marked endpoint and commuting trace
equivalence. A callback that inspects the particular certificate history
or arbitrarily selected witness cover is outside that guarantee. Strict
tetrahedron descent and the supplied-vector disc predicate have the required
invariance.

The finite-family cost separates three tasks: construction of the complete
cover index, exploration of geometric traces, and evaluation of endpoints.
The descent endpoint test is just a tetrahedron count. A normal-disc query
is optional and has its own potentially substantial cost. Performance
claims for the descent test therefore do not automatically transfer to an
arbitrary disc-search callback.

\subsection{Public interfaces and exact meanings of results}

The generic entry point is \code{search_pachner_cover} in
\code{fastunknot/pachner_cover_search.py}. It accepts a source triangulation
and height rows, a maximum region size, maximum number of components,
total upward budget, and explicit resource limits. The methods are the
maintained sleep-set route and a small-instance naive reference route.
The older restarted and commitment implementations remain available.

The convenience entry point \code{find_pachner_descent} supplies the
radius $\min(t,3U+3)$ unless the caller explicitly requests another radius.
It queries strict tetrahedron loss and independently verifies a positive
answer. A smaller requested radius is a legitimate local search, but its
failure has a deliberately narrower status.

\begin{table}[ht]
\centering\small
\begin{tabularx}{\textwidth}{@{}p{0.39\textwidth}X@{}}
\toprule
Status & Meaning \\
\midrule
\code{DESCENT_FOUND} & Source replay verifies strict descent and the declared
total upward bound. \\
\code{COMPLETE_BOUNDED_DESCENT} & No descent exists within that upward budget,
using the sufficient localization radius and complete search. \\
\code{COMPLETE_BOUNDED_LOCAL_DESCENT} & Only the explicitly smaller covered
family is exhausted. \\
\code{COMPLETE_BOUNDED_COVER_FAMILY} & The complete declared generic cover
family is exhausted. \\
\code{INCONCLUSIVE} & A required index, traversal, replay, or endpoint-query
operation reached its limit. \\
\bottomrule
\end{tabularx}
\caption{The principal search outcomes. Exhaustion is a bounded geometric
conclusion and is not a nontrivial-knot certificate.}
\label{tab:statuses}
\end{table}

The generic interface also distinguishes an independently verified normal-disc
answer, \code{DISC_FOUND}, from a caller-selected endpoint,
\code{ENDPOINT_SELECTED}. A certificate on an abstract manifold still needs
the maintained diagram-to-exterior binding before it can certify a statement
about a supplied knot diagram.

\begin{lstlisting}[language=Python,float=htbp,caption={Complete bounded descent with an independently checked positive witness.}]
from fastunknot.pachner_cover_search import find_pachner_descent
from fastunknot.pachner_cover_verify import verify_pachner_descent

answer = find_pachner_descent(
    triangulation, heights, max_upward=1,
    max_nodes=10000, max_regions=100000, max_work=2000000,
)
if answer["status"] == "DESCENT_FOUND":
    assert verify_pachner_descent(
        triangulation, heights, answer["certificate"],
        max_upward=1,
    )
\end{lstlisting}

\subsection{Independent certificate checking}

The producer emits the existing
\code{pachner-cochain-endpoint-v1} schema. Its declared active original
region is one concrete member of the compatible-cover set; it need not
equal the actual consumed footprint. The new verifier independently checks
that this cover has the stated size and number of induced dual components,
then invokes the maintained source-bound geometric and cochain replay.
It imports no region enumerator, search engine, native move producer, or
cochain-transport producer.

For descent, it also recomputes upward and downward counts, checks the
actual final tetrahedron count, and verifies
$t_{\rm final}=t_{\rm initial}+u-d<t_{\rm initial}$. This prevents the
search's status string or claimed counters from substituting for a proof.
The verification tests disable the corresponding producer functions and
still accept genuine retained certificates.

Resource caps are part of the observable contract. \code{max_regions}
limits construction of the \emph{complete} index, including the empty
region. Partial construction cannot establish complete exhaustion.
\code{max_nodes} counts visited search frames; \code{max_work} counts
cooperative checkpoints across preparation, indexing, search, and positive
replay. These checkpoints are operational accounting units, not bit
operations. Optional disc-query caps also prevent a complete disc-negative
claim when a required query was unfinished. User callback exceptions
propagate unchanged, including exceptions whose class happens to match an
internal budget exception. Callback inputs are isolated from mutable search
state.

\section{Certificate projection, composition, and packing}\label{sec:certificate-tools}

\subsection{Extracting a connected first descent}

The new \code{pachner_causality.py} reads a supplied endpoint proof only
after independent source replay accepts it. Each live tetrahedron receives
either a source identity or a pair consisting of its birth event and output
slot. Reading the consumed regions recovers every birth arc and the actual
source footprint. The analyzer reports the components, their losses, and
the incidence census.

For a descending certificate, \code{extract_local_descent} chooses a
positive-loss component, projects its events, and truncates at first
descent. These operations repeat until the retained trace is connected
and first-descending. One projection followed by one truncation is
insufficient: deleting a late event that joined earlier components can
disconnect the prefix again. Every necessary repetition strictly shortens
the trace, giving termination. The final witness satisfies the localization
bound and is independently replayed before return.

The consumer must respect every local frame accepted by the public verifier.
A valid downward certificate can use either order for two belt vertices in
the seed tetrahedron. A native producer may choose just one order. Replaying
with that producer's frame and then treating its born tetrahedra as if they
used the declared frame corrupts later dependent events, even though the
first replacement has the right unmarked geometry. The extractor therefore
retains the declared formal region and, where necessary, relabels the two
born tetrahedra and all incident permutations before carrying the cochain.
This fix is exercised by a retained noncanonical-frame counterexample and
an exhaustive finite frame audit described in \cref{sec:experiments}.

\subsection{Composing independently supplied traces}

The companion \code{compose_disjoint_traces} authenticates every input
endpoint against the same original source. It computes \emph{actual}
consumed footprints and rejects intersecting ones, even if their advertised
cover regions happen to differ. Birth names are namespaced by input trace,
so equal event numbers from different certificates cannot collide. The
replay transports source and birth identities, preserving declared formal
frames, and returns one ordinary source-bound endpoint certificate together
with the additive count census. This is an executable positive-witness
counterpart of \cref{lem:compose}. Overlapping supplied covers are harmless
when the actual footprints are disjoint.

The function does not assume that input traces are descending. The geometric
composition lemma applies to arbitrary legal traces with disjoint actual
footprints; a caller may separately inspect their losses. Empty traces are
allowed and contribute zero. Invalid proofs, overlap, or exhausted work
do not produce a partial successful composition.

\subsection{The supplied-coloring dynamic program}

The prototype \code{causal_research/colorful_packing.py} implements the
exact $0/1$ packing step of \cref{sec:multiunit} for a supplied coloring.
It validates nonempty distinct-source footprints, nonnegative upward costs,
and positive losses. It discards noncolorful items, compacts unused numeric
color labels, and preserves original candidate indices in immutable witness
tuples. A separate previous-item table prevents accidental reuse of an item.

Its objective is maximum capped loss, followed by minimum upward cost.
It does not claim globally minimum witness cardinality or lexicographically
first item indices: retaining only one best capped profit per mask and
cost does not justify those extra objectives. The result exposes actual
loss as well as capped loss, upward cost, union footprint, and color mask.
The candidate dictionaries can retain authenticated certificates, but the
combinatorial routine itself does not authenticate their geometry.

The delivered prototype does not generate the complete deterministic
perfect-hash family. Consequently a negative result from one coloring
is only a negative result for that colorful item instance. The mathematical
algorithm in \cref{thm:k-fpt}, the implemented packing primitive, and the
independent positive composition utility are separate, explicit deliverables.
